\documentclass[a4paper,12pt]{article}
\usepackage[utf8]{inputenc}
\usepackage[margin=0.9in]{geometry}
\usepackage{amsmath,amssymb}
\usepackage{caption}
\usepackage{verbatim}
\usepackage{booktabs}
\usepackage{hyperref}
\usepackage{microtype}
\sloppy

\newcommand{\aver}[1]{\ensuremath{\langle#1\rangle}}
\renewcommand{\t}{\ensuremath{\tau}}
\newcommand{\w}{\ensuremath{\omega}}
\newcommand{\TT}{\ensuremath{\mathcal{T}_\t}}
\newcommand{\Hloc}{\ensuremath{H_{\rm loc}}}
\newcommand{\ndag}{\ensuremath{^\dagger}}

\begin{document}
\title{Implementation notes: multi-orbital dynamical interactions and the\\
Lang--Firsov resummation in \texttt{CTHYB}}
%\author{Andrew Hardy}
\maketitle

\begin{abstract}
This note documents \texttt{CTHYB}'s multi-orbital retarded (dynamical) interaction
machinery: the 4-index vertex representation, the symbolic-algebra criterion that routes
each vertex to the analytic Lang--Firsov resummation or to the stochastic double
expansion, the linear-algebra extension of that criterion to \emph{combinations} of
orbital densities (which a genuine Hubbard--Kanamori $\Hloc$ requires), the static
chemical-potential and interaction shift the resummation induces, and the estimators that
read correlators off either sampling route. The static shift is derived in closed form
from the code's own boundary-value problem and cross-checked against the exactly-solvable
single-phonon polaron shift.
\end{abstract}

\tableofcontents

\section{Overview}

\begin{itemize}
  \item \textbf{The vertex.} A dynamical vertex is a pair of fermion bilinears with a
    retarded coupling, $D(\t)\,O_1(\t)\,O_2(0)$ (\texttt{dyn\_vertex\_t} in
    \texttt{configuration.hpp}), registered through
    \texttt{solver.add\_dyn\_vertex(op1,op2,coupling)}. \texttt{D0\_tau} and
    \texttt{Jperp\_tau} are convenience inputs that expand into the same representation
    (Sec.~\ref{sec:api}).
  \item \textbf{Routing.} A vertex is Lang--Firsov-eligible if both $O_1,O_2$ are number
    operators \emph{and} each commutes with $\Hloc$, checked by symbolic operator algebra
    in \texttt{classify\_dyn\_vertices} (Sec.~\ref{sec:eligibility}). Everything else is
    sampled stochastically.
  \item \textbf{Conserved density combinations.} A Kanamori $\Hloc$ with spin-flip and
    pair-hopping conserves no individual orbital density, only combinations of them
    ($N_\uparrow,N_\downarrow$ in the standard case).
    \texttt{find\_conserved\_density\_combinations} obtains that subspace as the nullspace
    of a commutator map, and \texttt{split\_density\_couplings} writes a rejected density
    coupling as an analytic part built only from those combinations plus a stochastic
    residual. The split is an identity, and is chosen so the residual carries no weight of
    the opposite sign (Secs.~\ref{sec:conserved}--\ref{sec:split-sign}).
  \item \textbf{The static shift.} The resummation is exact only once an instantaneous
    piece of the retarded coupling, $K'(0)$, is folded into $\Hloc$ before it is
    diagonalized. It renormalizes the effective interaction matrix and chemical potential,
    so a filling target computed from the bare interaction alone misses half filling and
    breaks particle-hole symmetry (Secs.~\ref{sec:shift}--\ref{sec:ph}).
  \item \textbf{Ergodicity of the stochastic expansion.} \texttt{insert\_dyn} and
    \texttt{remove\_dyn} act on whole vertices and cannot reach pairings in which every
    vertex crosses another. \texttt{swap\_dyn} exchanges partners between two vertices at
    fixed operator positions and makes the sampling ergodic (Sec.~\ref{sec:swap}).
  \item \textbf{Estimators.} Both routes admit an exact correlator estimator obtained by
    differentiating $\ln Z$ with respect to a retarded coupling: a vertex-separation
    histogram for the stochastic channels, and the density-kink correlator for the
    Lang--Firsov ones (Sec.~\ref{sec:derivative}).
\end{itemize}

\section{The vertex representation and the solve-time pipeline}
\label{sec:api}

\subsection{What a vertex is}

\texttt{dyn\_vertex\_t} (\texttt{configuration.hpp}) holds two
\texttt{many\_body\_operator} expressions and a coupling curve,
\begin{verbatim}
  struct dyn_vertex_t { many_body_op_t op1, op2; gf<imtime, scalar_valued> coupling; };
\end{verbatim}
meaning $D(\t)\,O_1(\t)\,O_2(0)$. The operators are written exactly as in
\texttt{h\_int}, e.g.\ \verb|c_dag('up',0)*c('down',0)|. This reduces to the block/inner-index
form \texttt{op\_desc\_pair\_t} by \texttt{extract\_bilinear} where the Monte Carlo needs
it. Three constraints follow from that lowering
\begin{itemize}
  \item \textbf{Each side is a single fermion bilinear} --- one creation and one
    annihilation operator. Pair hopping is therefore not representable:
    $c\ndag_a c\ndag_b$ is a pair-creation operator, and \texttt{insert\_dyn} /
    \texttt{remove\_dyn} assume one creation plus one annihilation per side. There is no
    \texttt{pair\_hopping} argument anywhere, so its absence cannot be mistaken for
    support.
  \item \textbf{Each side carries unit coefficient.} \texttt{extract\_bilinear} throws
    unless $|{\rm coeff}-1|<10^{-12}$. A scalar on $O_1$ or $O_2$ would be redundant with
    the coupling, $D\,(aM_1)(bM_2)=(abD)\,M_1M_2$, and absorbing it silently would hide a
    factor the caller can equally well write into $D$.
  \item \textbf{A sum of bilinears is expanded by the caller.} A spin operator in a
    rotated basis is a sum of several patch-basis bilinears, and has to be written out
    as monomial pairs; \texttt{expand\_retarded\_product} in the DCA benchmark's
    \texttt{dca\_model.py} does this and merges duplicates.
\end{itemize}

\subsection{The convenience inputs}

\texttt{D0\_tau} and \texttt{Jperp\_tau} expand into the same vertex list:
\begin{itemize}
  \item \texttt{expand\_D0\_into\_vertices}: every non-zero $(bl_1,i_1,bl_2,i_2)$ entry
    becomes its own vertex $n_{bl_1 i_1}(\t)n_{bl_2 i_2}(0)$. No orbital or spin
    convention is assumed, because a density-density coupling does not need one.
  \item \texttt{expand\_Jperp\_into\_vertices}: \texttt{Jperp\_tau} is a \emph{scalar}
    \texttt{gf<imtime>}, a single global up/down coupling matching CT-SEG's interface
    exactly. It therefore supports only the unambiguous case of exactly two blocks of
    one mode each, and raises an error otherwise, directing the caller to
    \texttt{add\_dyn\_vertex} --- anything more general would require guessing an orbital
    layout.
\end{itemize}
For a Kanamori model, \texttt{python/triqs\_cthyb/dynamical\_interactions.py} provides
\texttt{kanamori\_dynamical\_vertices(solver, spin\_names, orb\_names, U, Uprime, J\_hund,
spin\_flip)}, the dynamical generalization of \texttt{h\_int\_kanamori}'s density-density
and spin-flip terms term for term, with each coupling either a single \texttt{Gf}
broadcast over all pairs or a \texttt{dict[(a1,a2)] -> Gf}.

\subsection{The pipeline}

\texttt{solver\_core::solve()} runs, in order:
\begin{enumerate}
  \item \texttt{collect\_dyn\_vertices} --- merge the explicit vertices with the expanded
    \texttt{D0\_tau}/\texttt{Jperp\_tau} into one list.
  \item \texttt{classify\_dyn\_vertices} --- route each vertex to \texttt{.lang\_firsov}
    or \texttt{.stochastic} by the criterion of Sec.~\ref{sec:eligibility}.
  \item \texttt{conserved\_densities} $+$ \texttt{split\_density\_couplings} --- split the
    \emph{rejected density} vertices as in Sec.~\ref{sec:split}, moving part of each back
    to the analytic list.
  \item \texttt{apply\_lang\_firsov\_shift} --- fold the static $K'(0)$ part of the
    analytic vertices into $\Hloc$. This must run \emph{before} \texttt{h\_diag} is built.
  \item \texttt{build\_K\_n} --- build the $K_n[a][b][:]$ Legendre kernel used by
    \texttt{compute\_lang\_firsov\_ratio}.
  \item \texttt{fold\_into\_stochastic\_catalog} --- convert the rest into
    \texttt{dyn\_op\_list}/\texttt{dyn\_interactions} for \texttt{insert\_dyn} /
    \texttt{remove\_dyn} / \texttt{swap\_dyn}.
\end{enumerate}
Every vertex ends in exactly one list, so nothing is dropped, and
\texttt{has\_dyn\_interactions} --- which gates whether the stochastic moves and measure
are registered at all --- is simply \verb|!dyn_op_list.empty()|.

\subsection{Switches and the audit}

\texttt{params.lang\_firsov} defaults to \texttt{true}, which is safe because of the
per-vertex routing above: anything ineligible falls back to the stochastic expansion
automatically. Setting it to \texttt{false} skips the eligibility test entirely and forces
every vertex stochastic, which is what the \texttt{lang\_firsov=True}/\texttt{False}
cross-checks use.

At \texttt{verbosity >= 4} the solver prints an audit of everything it inferred: the
conserved density combinations of $\Hloc$; every vertex with the four conditions the
routing rests on (is it $n_a$, and does it commute with $\Hloc$, for each of $O_1,O_2$)
and the resulting route; the analytic and stochastic lists the split produced, with each
coupling's $D(0)$ and $D(\beta/2)$; the aggregate $K'(0)$ folded into $\Hloc$; and the
final $\Hloc$. (The compile-time \texttt{EXT\_DEBUG} option is separate: it prints per
Monte Carlo move, not once per solve.)

\section{Setup: retarded density-density interactions in the local trace}
\label{sec:setup}

Consider a local Hamiltonian $\Hloc$ (Kanamori interaction, crystal field, etc.) plus a
retarded density-density coupling between spin-orbitals $a,b$,
\begin{equation}
  S_{\rm ret} = \tfrac12 \int_0^\beta\!\!\int_0^\beta d\t\, d\t'\;
  n_a(\t)\, D_{ab}(\t-\t')\, n_b(\t'),
  \label{eq:sret}
\end{equation}
where $n_a(\t) = e^{\t \Hloc} n_a e^{-\t \Hloc}$ (Heisenberg picture under $\Hloc$
alone, between hybridization events) and $D_{ab}(\t)$ is whatever bosonic bath was
integrated out (a phonon, a plasmon, a $GW$-derived screened interaction, \ldots). In
this code $D_{ab}$ is exactly the \texttt{coupling} member of a
\texttt{dyn\_vertex\_t}. The concrete example used throughout
(\texttt{kanamori\_dyn.py}) is a single dispersionless phonon of frequency $\w_0$,
coupled with strength $g$ to every spin-orbital's density through the standard
harmonic-oscillator coordinate $x=(b+b\ndag)/\sqrt{2\w_0}$,
\begin{equation}
  H = \Hloc + \w_0 b\ndag b + g\Big(\sum_a n_a\Big) x .
  \label{eq:holstein}
\end{equation}
Integrating out the free phonon gives exactly the coupling used in the code,
\begin{equation}
  D(\t) = g^2\, Q(\t), \qquad
  Q(\t) = -\frac{1}{2\w_0}\,\frac{\cosh\!\big(\w_0(\t-\beta/2)\big)}{\sinh(\w_0\beta/2)},
  \qquad 0<\t<\beta,
  \label{eq:Qtau}
\end{equation}
identical for every pair $(a,b)$ since the coupling in Eq.~\eqref{eq:holstein} is
democratic. Two mechanisms can sample Eq.~\eqref{eq:sret}:
\begin{itemize}
  \item \textbf{Stochastic double expansion} (\texttt{moves/insert\_dyn.cpp},
    \texttt{remove\_dyn.cpp}): samples the double integral directly by inserting
    auxiliary vertices, always correct, always available.
  \item \textbf{Analytic Lang--Firsov resummation}: exact \emph{only} when $n_a$ (or a
    combination of densities) is a constant of motion between hybridization events;
    Secs.~\ref{sec:eligibility}--\ref{sec:conserved} make this precise.
\end{itemize}

\section{The Lang--Firsov resummation and its eligibility condition}
\label{sec:eligibility}

If $[n_a,\Hloc]=0$, then along the imaginary-time trace $n_a(\t)$ is a piecewise
constant, integer-valued function of $\t$ that only jumps when a $c_a\ndag$ or $c_a$
hybridization operator is inserted or removed. Under that condition the double
integral \eqref{eq:sret} can be performed exactly (Lang \& Firsov 1962
\cite{LangFirsov1962}; see \cite{WernerMillis2007,WernerMillis2010,GullRMP2011} for the
CT-HYB/CT-INT incarnation), and the local trace picks up an extra factor
per accepted/proposed move,
\begin{equation}
  w^{\rm dyn}_{\rm loc} = \exp\Big\{ \sum_{\alpha,\beta} s_\alpha s_\beta\,
  K_{a(\alpha)\,b(\beta)}(\t_\alpha-\t_\beta) \Big\},
  \label{eq:ratio}
\end{equation}
where the sum runs over pairs of fermion-operator insertions/removals $\alpha,\beta$
(each contributing $s=+1$ if it is a creation operator being inserted or an
annihilation operator being removed, $s=-1$ otherwise), and $K_{ab}$ is built once per
solve from $D_{ab}$ (\texttt{build\_K\_n}) and evaluated by
\texttt{qmc\_data::compute\_lang\_firsov\_ratio} (\texttt{qmc\_data.hpp}).

\texttt{classify\_dyn\_vertices} tests the condition by direct symbolic operator algebra,
$(O\,\Hloc-\Hloc\,O).\mathtt{is\_almost\_zero()}$, on the actual $\Hloc$ used in the
solve. It does not use \texttt{atom\_diag}, whose quantum-number machinery labels Fock
states by diagonal matrix elements in the \emph{fundamental} operator basis --- faithful
only when the candidate operator is already diagonal there --- and which segfaults during
construction when fed two or more simultaneously non-commuting \texttt{qn\_vector}
entries.

\section{From $D_{ab}(\t)$ to $K_{ab}(\t)$: a boundary-value problem in Legendre space}
\label{sec:shift}

\subsection{The pairwise kernel}

\texttt{fit\_legendre\_coeffs} expands the coupling curve in ordinary (non-normalized)
Legendre polynomials on $\t\in[0,\beta]$,
\begin{equation}
  D(\t) = \sum_{n\ge0} d_n\, P_n\big(x(\t)\big), \qquad x(\t) = \frac{2\t}{\beta}-1,
  \qquad d_n = \frac{2n+1}{\beta}\int_0^\beta \! d\t\, D(\t)\, P_n(x(\t)).
  \label{eq:dn}
\end{equation}
\texttt{build\_M\_matrix} then maps $\{d_n\}\to\{k_n\}$, the Legendre coefficients of
$K(\t)=\sum_n k_n P_n(x(\t))$, via a fixed matrix independent of $D$. One can check
directly that this matrix is exactly the representation, in Legendre space, of the
two-point boundary-value problem
\begin{equation}
  K''(\t) = D(\t), \qquad K(0)=K(\beta)=0,
  \label{eq:bvp}
\end{equation}
i.e.\ $K$ is the doubly-integrated coupling curve, pinned to vanish at both ends of the
trace. (Check: for $D(\t)=d_0$ constant, Eq.~\eqref{eq:bvp} gives
$K(\t)=\tfrac{d_0}{2}\t(\t-\beta) = \tfrac{d_0\beta^2}{12}\big(P_2(x)-P_0(x)\big)$ using
$x^2=\tfrac13(2P_2+P_0)$, i.e.\ $k_0=-d_0\beta^2/12$, $k_2=+d_0\beta^2/12$ ---
exactly \texttt{build\_M\_matrix}'s $M(0,0)=-\beta^2/12$ and $M(2,0)=+\beta^2/12$.)
Physically, $K$ is the retarded pair potential between two density kinks a time $\t$
apart; the Dirichlet condition $K(0)=0$ is why \texttt{compute\_lang\_firsov\_ratio}'s
comment notes that coincident-time pairs contribute nothing to Eq.~\eqref{eq:ratio}.

\subsection{The static part, $K'(0)$}
\label{sec:kprime}

What is \emph{not} captured by pairs of distinct kinks is the instantaneous
(Hartree-like) part of the interaction --- exactly the piece that must instead be
folded into $\Hloc$ before it is diagonalized (\texttt{apply\_lang\_firsov\_shift}), or
it is missed entirely. Solving the BVP \eqref{eq:bvp} with the Dirichlet Green's
function of $\partial_\t^2$ on $[0,\beta]$,
$G(\t,\t') = \t_<(\t_>-\beta)/\beta$ ($\t_<=\min(\t,\t')$, $\t_>=\max(\t,\t')$), gives
\begin{equation}
  K'(0^+) = \int_0^\beta \frac{\t'-\beta}{\beta}\, D(\t')\, d\t'
  = \frac1\beta\int_0^\beta \t' D(\t')\,d\t' \;-\; \int_0^\beta D(\t')\,d\t'.
\end{equation}
Both integrals are exactly the $n=0,1$ Legendre moments of $D$ (using
$x^0=P_0$, $x=P_1$, and orthogonality $\int_{-1}^1 P_nP_m\,dx = \tfrac{2}{2n+1}\delta_{nm}$):
$\int_0^\beta D\,d\t = \beta d_0$ and $\int_0^\beta \t' D\,d\t' =
\tfrac{\beta^2}{2}\big(d_0+\tfrac{d_1}{3}\big)$, giving the closed form
\begin{equation}
  K'(0^+) = -\frac{\beta}{2}\Big(d_0 - \frac{d_1}{3}\Big).
  \label{eq:kprime0plus}
\end{equation}
This uses \emph{only} the first two Legendre moments of $D$ --- exact, not a
truncation. The kernel \texttt{compute\_lang\_firsov\_ratio} actually evaluates is not
$K(\t)$ on $[0,\beta]$ as-is, but folded/reflected into $[0,\beta/2]$
(\texttt{eval\_K}'s \texttt{if (t > beta/2.0) t = beta - t;} after wrapping negative
$\t$ into $[0,\beta)$), which is equivalent to extending $K$ \emph{evenly} about
$\t=0$. Since $K'(0^+)\neq0$ in general, this even extension has a genuine kink at the
origin: approaching from below, $K_{\rm used}(\t)=K(-\t)$ gives
$K_{\rm used}'(0^-) = -K'(0^+)$, so the total derivative discontinuity at coincidence
is
\begin{equation}
  K'(0) \;\equiv\; K_{\rm used}'(0^+) - K_{\rm used}'(0^-) \;=\; 2K'(0^+)
  \;=\; -\beta\Big(d_0-\frac{d_1}{3}\Big),
  \label{eq:kprime0}
\end{equation}
which is exactly \texttt{apply\_lang\_firsov\_shift}'s
\verb|Kprime_0 = -beta*(d0 - d1/3)|. Physically, this discontinuity is the
correction a single density segment's own self-interaction contributes as its two
bounding kinks are brought together --- the Hartree/instantaneous piece the pairwise
sum in Eq.~\eqref{eq:ratio} genuinely cannot see (it starts at $K(0)=0$, not at the
one-sided slope). \texttt{apply\_lang\_firsov\_shift} moves it into $\Hloc$ instead:
\begin{align}
  a=b\ (\text{diagonal}): &\quad \Hloc \;\to\; \Hloc - \tfrac12 K'(0)\, n_a, \\
  a\neq b\ (\text{off-diagonal}): &\quad \Hloc \;\to\; \Hloc - \tfrac12 K'(0)\, n_a n_b
\end{align}
applied for both orderings $(a,b),(b,a)$, giving a total $K'(0)$ per unordered pair.

\subsection{Computing the shift in practice}
\label{sec:kprime-practice}

Four facts about Eq.~\eqref{eq:kprime0} that are easy to get wrong:

\begin{itemize}
  \item \textbf{The two-term formula is an identity, not a truncation.} Equivalently
    $K'(0)=-\tfrac1\beta\int_0^\beta(\beta-\t)D(\t)\,d\t$, and $(\beta-\t)$ is degree 1
    in $x=2\t/\beta-1$, so every $P_{l\ge2}$ is orthogonal to it and cannot contribute.
    \textbf{\texttt{dyn\_n\_l} therefore cannot change the static offset}, and raising it
    cannot fix an offset problem. The only error is the quadrature of the two moments:
    the C++ \texttt{fit\_legendre\_coeffs} uses the trapezoid rule, whose error is
    ${\sim}10^{-6}$ at \texttt{n\_tau\_bosonic}${}=2001$ because the boson kernel is most
    curved exactly at the endpoints; the Python \texttt{kprime\_0} uses Simpson
    (${\sim}10^{-10}$).
  \item \textbf{Closed form for a boson.} $\int_0^\beta Q = -1/\w_0^2$ and
    $\int_0^\beta Q\,x = 0$, so a coupling $\lambda\,Q(\t)$ has
    $K'(0)=\lambda/(2\w_0^2)$ exactly. For the phonon $D_{ab}=g_ag_bQ$ this is an
    unordered-pair shift $-g_ag_b/\w_0^2$ and a level shift $-g_a^2/(2\w_0^2)$, which is
    what the cross-check below rederives.
  \item \textbf{The zero-frequency shortcut needs a symmetric kernel.} If
    $D(\t)=D(\beta-\t)$ then $\langle\beta-\t\rangle=\beta/2$ and
    $K'(0)=-\tfrac12 D(i\nu=0)$, as \texttt{holstein.py} uses. Ordinary boson propagators
    are symmetric, so it is correct there; for an asymmetric kernel (a charged or complex
    boson, or a $D(\t)$ coming out of a self-consistent loop) it is wrong --- by 44\% for
    the control case in the \texttt{Py\_dyn\_static\_shift} test.
  \item \textbf{Sign trap.} The level shift enters
    $\mu_a=\tfrac12\sum_{b\neq a}W_{ab}+{\rm level\_shift}_a$ with a \emph{plus}, because
    it is the shift applied to the orbital \emph{energy} ($=-K'(0)$), not to $\mu$. With
    the sign flipped, the Kanamori-plus-phonon case $g=(0.7,0.3)$ returns a $\mu$ that is
    \emph{uniform} across orbitals ($2.04$ for both) instead of $1.55/1.95$ --- entirely
    plausible-looking and wrong. The test asserts that the two orbitals differ for
    exactly this reason.
\end{itemize}

\subsection{Cross-check: the exactly-solvable single-phonon polaron shift}

For the concrete model \eqref{eq:holstein}, the phonon coordinate can be integrated
out exactly (this is the original Lang--Firsov canonical transformation, and also just
completing the square classically, since the coupling is linear): minimizing
$g\,n\,x+\tfrac{\w_0^2}{2}x^2$ over $x$ at fixed $n=\sum_a n_a$ gives
$x^\star=-gn/\w_0^2$ and an energy shift $-\tfrac{g^2}{2\w_0^2}n^2$. Writing
$n^2=\sum_a n_a + \sum_{a\neq b} n_an_b$ (fermion idempotency, $n_a^2=n_a$) gives an
exact, textbook chemical-potential shift of $-\tfrac{g^2}{2\w_0^2}$ per orbital and an
exact interaction shift of $-\tfrac{g^2}{\w_0^2}$ per unordered pair. Independently,
\eqref{eq:Qtau}--\eqref{eq:kprime0} give, using
$\int_0^\beta Q(\t)\,d\t = -1/\w_0^2$ (elementary, and $\beta$-independent) so that
$d_0 = g^2\!\int Q\,d\t/\beta = -g^2/(\w_0^2\beta)$ and $d_1=0$ (by the manifest
symmetry $Q(\t)=Q(\beta-\t)$),
\begin{equation}
  K'(0) = -\beta\Big({-\frac{g^2}{\w_0^2\beta}}\Big) = \frac{g^2}{\w_0^2}.
\end{equation}
Hence $-\tfrac12 K'(0) = -\tfrac{g^2}{2\w_0^2}$ (diagonal) and $-K'(0) =
-\tfrac{g^2}{\w_0^2}$ (off-diagonal),
identical to the classical/exact Lang--Firsov result above, and independent of
$\beta$ as it must be (a $T=0$-like intrinsic property of the phonon; only the
pairwise, retarded part $K(\t)$ in Sec.~\ref{sec:shift} knows about $\beta$). For
$g=0.5,\ \w_0=1$: $K'(0) = 0.25$.

\section{Consequence: particle-hole symmetry and the $\mu$ correction}
\label{sec:ph}

For a purely static density-density Hamiltonian $\sum_{i<j}U_{ij}n_in_j$, half filling
($\langle n_i\rangle=\tfrac12$ for every spin-orbital $i$) is obtained at
$\mu_i = \tfrac12\sum_{j\neq i}U_{ij}$. \texttt{kanamori\_dyn.py} computes exactly this
for the bare Kanamori $H_{\rm int}$ alone,
\begin{equation}
  \mu_i^{\rm bare} = \tfrac12 U + (n_{\rm orb}-1)\tfrac12(U-2J) + (n_{\rm orb}-1)\tfrac12(U-3J),
\end{equation}
which is correct for $H_{\rm int}$ by itself (spin-flip and pair-hopping are
off-diagonal 4-fermion processes and do not shift this formula --- see
Sec.~\ref{sec:kanamori-check} for why they conserve $N_\uparrow,N_\downarrow$
individually, which is what makes the standard PH transformation still apply). But
once the phonon's static part (Sec.~\ref{sec:kprime}) is folded into $\Hloc$, the
\emph{total} interaction matrix seen by the model is $U^{\rm renorm}_{ij} =
U^{\rm bare}_{ij} - K'(0)_{ij}$, and the true half-filling point is
$\mu_i^{\rm half} = \tfrac12\sum_{j\neq i}U^{\rm renorm}_{ij}$, not $\mu_i^{\rm bare}$.
Since \texttt{solve()} \emph{also} applies the diagonal shift internally
($\mu_i^{\rm eff} = \mu_i^{\rm input} + \tfrac12 K'(0)_{ii}$ automatically, whatever
input $\mu$ is supplied), the input one must actually pass is
\begin{equation}
  \mu_i^{\rm input} \;=\; \mu_i^{\rm half} - \tfrac12 K'(0)_{ii},
  \label{eq:mucorr}
\end{equation}
so that the two shifts compose to the correct target.

\paragraph{Who computes it.} Four pure functions in
\texttt{python/triqs\_cthyb/dynamical\_interactions.py} --- \texttt{kprime\_0},
\texttt{kprime\_0\_boson}, \texttt{static\_shift} and \texttt{half\_filling\_mu} ---
evaluate Eqs.~\eqref{eq:kprime0} and \eqref{eq:mucorr} from the \emph{input} vertex list,
before any solve, and the caller adds the result to $\mu$ explicitly. The input coupling
is the right starting point because $K'(0)$ is a property of the Hamiltonian, not of the
route that samples it: a non-uniform coupling has part of its static shift inside vertices
that go stochastic, so any quantity reported by the analytic path alone under-counts it,
and a $\mu$ taken from such a quantity would give \texttt{lang\_firsov=True} and
\texttt{lang\_firsov=False} runs \emph{different Hamiltonians}, destroying the
cross-check the two routes exist to provide. For the same reason a benchmark driver that
reads the shift back from a probe solve must run that probe with
\texttt{lang\_firsov=True} whatever the production setting, and save the resulting $\mu$
so the analysis can confirm both runs used it. For the realized Hamiltonian itself,
\texttt{h\_loc()} returns it exactly, including non-density terms that a $U/\mu$ pair
cannot represent. \texttt{Py\_dyn\_static\_shift} pins the four functions against
\texttt{ed\_reference}'s ED-validated $\mu$ and against \texttt{holstein.py}'s
$U/2-g^2/\w_0^2$.

\subsection{What it looks like when this is skipped}

For the two-orbital benchmark ($U=2.0$, $J=0.2$, $g=0.5$, $\w_0=1$, $\beta=10$, so
$K'(0)=0.25$ per Sec.~\ref{sec:shift}, uniform over all $2n_{\rm orb}=4$
spin-orbitals): $\mu^{\rm bare}=2.5$, and Eq.~\eqref{eq:mucorr} with 3 other
spin-orbitals per site gives
\begin{equation}
  \mu^{\rm half} = 2.5 - \tfrac32(0.25) = 2.125, \qquad
  \mu^{\rm input} = 2.125 - \tfrac12(0.25) = 2.0,
\end{equation}
matching the actual solver output of the corrected run
to 5 significant figures. A production run that predates this correction (same
physical parameters, run before the fix landed) used the uncorrected
$\mu^{\rm bare}=2.5$ throughout, and shows exactly the expected symptom in the raw
data:
\begin{center}
\begin{tabular}{lccc}
  \toprule
  & $G(0)$ & $G(\beta)$ & $\max_\t|G(\t)-G(\beta-\t)|$ \\
  \midrule
  uncorrected ($\mu=2.5$) & $-0.469$ & $-0.531$ & up to $0.55$ \\
  corrected ($\mu=2.0$)   & $-0.499$ & $-0.501$ & $\le 0.003$ \\
  \bottomrule
\end{tabular}
\end{center}
i.e.\ the uncorrected run sits at filling $n\approx0.47$ rather than $0.5$ per
spin-orbital --- precisely the density shift implied by a $0.5$ error in $\mu$ against
an interaction scale $U\sim2$ --- and the corrected run recovers $\t\to\beta-\t$
symmetry to noise level, matching a $g=0$ control run solved with the same code.

\section{Beyond a single density: conserved combinations}
\label{sec:conserved}

Section~\ref{sec:eligibility}'s criterion is correct but conservative: a real Kanamori
$\Hloc$ with spin-flip and pair-hopping does not conserve any individual $n_a$, only
certain \emph{combinations}.

\subsection{Why $N_\uparrow,N_\downarrow$ are conserved but individual $n_a$ are not}
\label{sec:kanamori-check}

The Kanamori spin-flip term is $H_{\rm sf} = -J_X\sum_{m\neq m'} d_{m\uparrow}\ndag
d_{m\downarrow}^{\phantom\dagger} d_{m'\downarrow}\ndag d_{m'\uparrow}^{\phantom\dagger}$.
Writing it as $d_{m\uparrow}\ndag\, B\, d_{m'\uparrow}$ with
$B=d_{m\downarrow}d_{m'\downarrow}\ndag$ (which commutes with $N_\uparrow$, since $B$
only involves down operators),
\begin{equation}
  [N_\uparrow,\, d_{m\uparrow}\ndag B d_{m'\uparrow}]
  = (+1)\,d_{m\uparrow}\ndag B d_{m'\uparrow} + d_{m\uparrow}\ndag B\,(-1)\,d_{m'\uparrow} = 0
\end{equation}
exactly, for every $J_X$: an up-electron hops $m'\!\to\!m$ while a down-electron hops
$m\!\to\!m'$ \emph{simultaneously}; no single electron ever changes spin, so
$N_\uparrow$ (and by the identical argument $N_\downarrow$) is untouched. The
pair-hopping term $J_P\sum_{m\neq m'}d_{m\uparrow}\ndag d_{m\downarrow}\ndag
d_{m'\downarrow}d_{m'\uparrow}$ moves one up- and one down-electron together from
$m'\to m$, again leaving $N_\uparrow,N_\downarrow$ individually fixed. Neither term
conserves any single orbital's $n_{m\uparrow}$ or $n_{m\downarrow}$.

\subsection{The conserved subspace as a nullspace}

$[\,\sum_a c_a n_a,\ \Hloc\,]=0$ is linear in $c$, so the set of conserved
combinations is exactly the nullspace of the linear map $c \mapsto \sum_a c_a\,[n_a,\Hloc]$.
\texttt{find\_conserved\_density\_combinations} builds this map by expanding every
$[n_a,\Hloc]$ symbolically, collecting the (real and imaginary parts of the)
coefficient of every monomial that appears as a row of a real matrix, and computing
its nullspace by SVD --- for the standard 2-orbital Kanamori case this returns
exactly the 2-dimensional space spanned by $N_\uparrow,N_\downarrow$, matching
Sec.~\ref{sec:kanamori-check} directly rather than assuming it.

A vertex whose two densities individually fail Sec.~\ref{sec:eligibility} is neither
simply accepted nor simply rejected: its coupling is split (Sec.~\ref{sec:split}) into a
part built only from the conserved combinations, $\sum_{ij}\Gamma_{ij}\,O_i(a)O_j(b)$,
which is resummed analytically, and a residual, which is sampled stochastically. A
coupling that is entirely of that form leaves no residual and goes analytic in full.
Because $N_\uparrow$ and $N_\downarrow$ have disjoint orbital support, same-spin and
opposite-spin pairs land on non-overlapping matrix entries, so $U(\t)\neq U'(\t)$ (the
physically standard Kanamori shape, not a degenerate special case) is of that form and
is handled entirely analytically
(\texttt{test/python/kanamori\_dyn\_nonuniform.py}: 16/16 vertices analytic).

\section{Splitting a density coupling into Lang--Firsov and stochastic parts}
\label{sec:split}

Implemented as \texttt{split\_density\_couplings}
(\texttt{dynamical\_interactions.cpp}), with the sign-preserving choice of
Sec.~\ref{sec:split-sign}.

\subsection{The coupling as a matrix, and the projector $P$}

Collect the density-density vertices into an $M\times M$ matrix $D(\t)$ over the $M$
spin-orbitals, with $D_{ab}=0$ for every pair the user did not specify, and write
$\mathbf n=(n_1,\dots,n_M)^T$, so that the density part of Eq.~\eqref{eq:sret} reads
$S_{\rm ret}=\tfrac12\int\!\!\int d\t\,d\t'\,\mathbf n(\t)^T D(\t-\t')\,\mathbf n(\t')$.
Vertices that are not density-density (spin-flip, \ldots) are not part of $D$ and are
untouched by what follows.

A density combination $\sum_a c_a n_a = c\cdot\mathbf n$ is labelled by its coefficient
vector $c\in\mathbb R^M$. The conserved ones form the subspace
$\mathcal C=\{c : [c\cdot\mathbf n,\Hloc]=0\}$ of Sec.~\ref{sec:conserved}. Pick an
orthonormal basis $v_1,\dots,v_r$ of $\mathcal C$ (the SVD null vectors returned by
\texttt{find\_conserved\_density\_combinations} already are one) and let
$O_i=v_i\cdot\mathbf n$. Define
\begin{equation}
  P=\sum_{i=1}^r v_i v_i^T ,
  \label{eq:projector}
\end{equation}
the orthogonal projector onto $\mathcal C$: a real symmetric $M\times M$ matrix with
$P^2=P$, $Pc=c$ for $c\in\mathcal C$ and $Pc=0$ for $c\perp\mathcal C$. It depends only
on $\mathcal C$, not on the basis chosen (for a non-orthonormal basis stacked as the rows
of $W$, $P=W^T(WW^T)^{-1}W$), and only on $\Hloc$, not on $\t$ or on the Monte Carlo
configuration. $P$ acts on coefficient vectors, not on Fock space. Acting on the vector
of densities, $(P\mathbf n)_a=\sum_i v_i[a]\,O_i$ is the conserved part of $n_a$: the
piece of it the Lang--Firsov resummation can see.

\subsection{The split and why it is exact}

Write
\begin{equation}
  D(\t)=P\,D(\t)\,P+R(\t),\qquad R\equiv D-PDP .
  \label{eq:split}
\end{equation}
This is an identity, and $S_{\rm ret}$ is linear in $D$, so treating $PDP$ and $R$ by
different methods changes no physics and adds no coupling the user did not ask for.

\paragraph{$PDP$ couples only conserved combinations.}
$\mathbf n^T PDP\,\mathbf n=(P\mathbf n)^T D\,(P\mathbf n)=\sum_{ij}\Gamma_{ij}(\t)\,O_iO_j$
with $\Gamma_{ij}=v_i^TDv_j$. Each $O_i(\t)$ is piecewise constant, jumping by
$\pm v_i[a]$ at every creation/annihilation operator on orbital $a$ in the trace
(hybridization and stochastic-vertex operators alike, cf.\ \texttt{qmc\_data::trace\_ops}),
so the argument of Sec.~\ref{sec:eligibility} applies unchanged. The code never needs
$\Gamma$ explicitly: $K$ and $K'(0)$ are linear in the coupling, so building the
orbital-resolved $K_{ab}$ of Eq.~\eqref{eq:ratio} and the static shift directly from
$(PDP)_{ab}=\sum_{ij}v_i[a]v_j[b]\Gamma_{ij}$ gives
$\sum_{\alpha\beta}s_\alpha s_\beta K_{a_\alpha b_\beta}
=\sum_{ij}\sum_{\alpha\beta}(v_i[a_\alpha]s_\alpha)(v_j[b_\beta]s_\beta)K^{\Gamma}_{ij}$,
the Lang--Firsov weight of $\sum_{ij}\Gamma_{ij}O_iO_j$, together with a static shift
$-\tfrac12\sum_{ij}K'^{\Gamma}_{ij}(0)\,O_iO_j$ that commutes with $\Hloc$. Being a
polynomial in densities, the shift also leaves $\mathcal C$ unchanged.

\paragraph{$R$ has no conserved--conserved part.}
Since $Pv_i=v_i$, $v_i^TRv_j=v_i^TDv_j-v_i^TDv_j=0$ for all $i,j$: every term left in
$R$ has at least one non-conserved factor, so it cannot be resummed and goes to the
stochastic expansion, which is exact for any coupling.

\paragraph{The Lang--Firsov part is as large as possible.}
A density coupling can be resummed only if it is of the form $\sum_{ij}\Gamma_{ij}O_iO_j$,
i.e.\ lies in $\mathcal L=\{E:PEP=E\}$. $PDP$ is the orthogonal (Frobenius) projection of
$D$ onto $\mathcal L$ ($PRP=0$ implies ${\rm tr}(R^TE)=0$ for every $E\in\mathcal L$), so
$R$ is as small as it can be, and $R=0$ exactly when $D$ is representable
($D=V^T\Gamma V\Leftrightarrow PDP=D$, with $V$ the matrix of rows $v_i$), which is
exactly the condition under which the whole coupling is analytic. Smallest is not the
same as best:
Sec.~\ref{sec:split-sign} explains why the implementation uses a different, sign-preserving
choice.

\paragraph{Practicalities.} $P$ is $\t$-independent, so Eq.~\eqref{eq:split} is applied
pointwise on the $\t$ grid, with no Legendre refit. $P$ is symmetric, so if
$D_{ab}(\t)=D_{ba}(\beta-\t)$ both pieces inherit this symmetry and are valid couplings
in the same vertex convention as $D$.

\subsection{Examples}

\begin{itemize}
  \item \textbf{Density-only $\Hloc$.} Every $n_a$ is conserved, $\mathcal C=\mathbb R^M$,
    $P=1$, $R=0$: the whole coupling is resummed analytically.
  \item \textbf{Only $N$ conserved.} $v=\mathbf 1/\sqrt M$, $P=\mathbf 1\mathbf 1^T/M$ and
    $PDP=\bar U(\t)\,\mathbf 1\mathbf 1^T$ with $\bar U=M^{-2}\sum_{ab}D_{ab}$: exactly the
    ``$U_{\rm tot}+\delta U$'' split, with $U_{\rm tot}$ the average of all entries and
    $\delta U_{ab}=D_{ab}-\bar U$.
  \item \textbf{Kanamori, two orbitals} ($N_\uparrow,N_\downarrow$ conserved, ordering
    $0\!\uparrow,1\!\uparrow,0\!\downarrow,1\!\downarrow$):
    \[
      P=\frac12\begin{pmatrix}1&1&0&0\\1&1&0&0\\0&0&1&1\\0&0&1&1\end{pmatrix},
    \]
    so $PDP$ replaces every entry of each spin block $(\sigma,\sigma')$ of $D$ by that
    block's average. A phonon on orbital 0 only, $D=g^2Q(\t)\,uu^T$ with $u=(1,0,1,0)$,
    gives $Pu=\tfrac12\mathbf 1$, $PDP=\tfrac14g^2Q\,\mathbf 1\mathbf 1^T$ and
    $R=g^2Q\,(uu^T-\tfrac14\mathbf 1\mathbf 1^T)$. Equivalently, with
    $n_0=\tfrac12N+\tfrac12\delta$ and $\delta=n_0-n_1$ the orbital polarization,
    $n_0n_0=\tfrac14NN+\tfrac14(N\delta+\delta N)+\tfrac14\delta\delta$: Lang--Firsov takes
    $\tfrac14NN$, the stochastic expansion the rest.
\end{itemize}

\paragraph{$P$ uses the full $v_i$.} Conservation is a property of $\Hloc$ on the whole
orbital space, not of the orbitals one vertex happens to touch, so the $v_i$ are never
restricted to a vertex's own support before projecting. Restricting them in the last
example would leave $n_{0\uparrow}$ and $n_{0\downarrow}$, making the coupling trivially
representable and routing all four orbital-0 vertices to Lang--Firsov --- treating the
non-conserved $n_0$ as conserved. With the full $v_i$, $Pu\neq u$ and the residual
survives.

\subsection{Which split: the sign}
\label{sec:split-sign}

Any Lang--Firsov part $D_{\rm LF}\in\mathcal L$ gives an exact split $D=D_{\rm LF}+R$: the
argument for $PDP$ above only used $D_{\rm LF}\in\mathcal L$. $PDP$ makes $R$ as small as
possible, but $R$ can then take either sign, and can be non-zero where $D$ is zero, while a
stochastic vertex carries weight $-R_{ab}(\t)$. Take the two-orbital Kanamori $\Hloc$ with a
phonon of orbital-dependent coupling, $D_{ab}=g_ag_bQ$ with $(g_0,g_1)=(0.7,0.3)$. Every entry of
$D$ has the sign of $Q<0$, so the fully stochastic expansion has only positive weights. $PDP$
replaces each spin block by its average, $0.25\,Q$, and leaves
$R/Q = 0.49-0.25 = +0.24$, $0.21-0.25 = -0.04$, $0.09-0.25 = -0.16$ for orbital pairs $00$,
$01$, $11$: two of the three residual couplings have the opposite sign. The measured average
sign drops from $1.0$ (fully stochastic) to $0.23$ ($\beta=5$, $U=2$, $J=0.3$, $V=0.7$,
$\w_0=1$). Both agree with exact diagonalization, but the statistical error grows as
$1/\text{sign}$.

\texttt{split\_density\_couplings} therefore uses a sign-preserving split. When the conserved
combinations are indicator vectors of disjoint sets of orbitals $S_i$ (as for $N$, or
$N_\uparrow,N_\downarrow$; \texttt{conserved\_densities} returns them in this form), $\mathcal L$
is the set of matrices that are constant on each block $S_i\times S_j$ and zero outside, and for
each $\t$ and each block the Lang--Firsov part is
\begin{equation}
  c_{ij}(\t)=\begin{cases}
    \min_{a\in S_i,\,b\in S_j} D_{ab}(\t) & \text{if every entry of the block is} >0,\\
    \max_{a\in S_i,\,b\in S_j} D_{ab}(\t) & \text{if every entry is} <0,\\
    0 & \text{otherwise,}
  \end{cases}
  \label{eq:signsplit}
\end{equation}
the entry of smallest magnitude when all entries share a sign. Every entry of $R=D-c$ then has
the sign of the corresponding entry of $D$ or vanishes, so the split adds no weight of the opposite
sign; for a block of one sign, $c$ is the largest block-constant part with that property. In the
example, $c=0.09\,Q$ and $R/Q=0.40,\ 0.12,\ 0$. An exactly representable coupling is
block-constant, so $c=D$ and $R=0$, as with $PDP$. If some pair in a block is not coupled at all,
$c=0$ for that block: the phonon on orbital 0 only stays fully stochastic, which is sign-free
there. $c(\t)$ is continuous but has kinks where the minimizing pair changes with $\t$; it is
resummed through its Legendre expansion like any coupling, so \texttt{dyn\_n\_l} must resolve
it. If the conserved combinations are not of the disjoint form, the rejected density vertices
stay fully stochastic; a general sign-preserving split would be a small linear program per $\t$,
not implemented.

\section{Measuring correlators as coupling derivatives}
\label{sec:derivative}

Every correlator in this section is a derivative of $\ln Z$ with respect to a retarded
coupling, estimated in whichever representation carries that coupling. This is the same
logic as the standard Green's function estimator,
$G(\t)=-\tfrac1\beta\aver{\sum_{ij}M^{-1}_{ji}\,\delta(\t-(\t_i-\t_j))}
=\delta\ln Z/\delta\Delta$~\cite{GullRMP2011}, and it fails in the same way: the
estimator divides by the coupling, so it becomes noisy as the coupling goes to zero.
There are two cases:
\begin{itemize}
  \item a coupling direction handled by Lang--Firsov gives the kink estimator
    (\texttt{measure\_D0\_corr}), exact even at zero coupling
    (Sec.~\ref{sec:kink-derivative});
  \item a coupling sampled stochastically gives the vertex-separation histogram
    (Sec.~\ref{sec:histogram}), exact but noisy wherever the coupling is small.
\end{itemize}
Directions that are neither (a non-conserved density pair with no stochastic vertex)
need an insertion estimator, which is not covered here.

Throughout, $\aver{X}$ is the Monte Carlo average including the sign $s$ and
\texttt{atomic\_reweighting} $r$: $\aver{X}=\sum sr\,X/\sum sr$. We assume the coupling
is real and symmetric,
\begin{equation}
  D_{ab}(\t)=D_{ba}(\t)=D_{ab}(\beta-\t),
  \label{eq:dsym}
\end{equation}
as for any coupling obtained from a bosonic bath with real vertices (the phonon
$g_ag_bQ(\t)$ of Eq.~\eqref{eq:Qtau}, or a cRPA/$GW$ $U_{ab}(\t)$ in a real orbital
basis). \texttt{eval\_K}'s folding of $\t$ into $[0,\beta/2]$ already relies on this.

\subsection{What the stochastic moves sample}
\label{sec:sampled-weight}

\texttt{insert\_dyn} draws $\t_1,\t_2$ uniformly and orders them so that $\t_1>\t_2$. It
picks a catalogue entry (type) $t$ uniformly among $N_t$, which gives a proposal density
$2/(\beta^2N_t)$, and has weight ratio $-f_t(\t_1-\t_2)$ times the trace and Lang--Firsov
ratios. \texttt{remove\_dyn} picks one of the $k$ vertices uniformly and uses the
inverse ratio. Detailed balance then fixes the stationary weight of $k$ vertices as
\[
  \frac1{k!}\prod_{v}\big[-f_{t_v}(s_v)\,d\t_{1v}\,d\t_{2v}\big]\times
  \mathrm{Tr}\times\det\times w^{\rm dyn}_{\rm loc},\qquad s_v=\t_{1v}-\t_{2v}\in(0,\beta),
\]
the time-ordered expansion of $Z=\mathrm{Tr}\,\TT\,e^{-S_{\rm samp}}$ with
\begin{equation}
  S_{\rm samp}=\sum_t\int\!\!\int_{\t_1>\t_2}d\t_1\,d\t_2\;f_t(\t_1-\t_2)\,
  O_{1,t}(\t_1)\,O_{2,t}(\t_2).
  \label{eq:ssamp}
\end{equation}
Relation to Eq.~\eqref{eq:sret}: split the square $[0,\beta]^2$ into $\t>\t'$ and
$\t<\t'$. In the second region use bosonic periodicity, $D(s-\beta)=D(s)$, and rename
$\t\leftrightarrow\t'$, $a\leftrightarrow b$ (densities commute under $\TT$):
\[
  S_{\rm ret}=\sum_{ab}\int\!\!\int_{\t>\t'}n_a(\t)\,
  \tfrac12\big[D_{ab}(s)+D_{ba}(\beta-s)\big]\,n_b(\t'),\qquad s=\t-\t'.
\]
Under Eq.~\eqref{eq:dsym} the bracket is $D_{ab}(s)$. So every ordered pair $(a,b)$ is
its own type with $f_{(a,b)}=D_{ab}$ at full weight, which is exactly how
\texttt{expand\_D0\_into\_vertices} builds the catalogue.

\subsection{Re-pairing the vertices: \texttt{swap\_dyn}}
\label{sec:swap}

\texttt{insert\_dyn} and \texttt{remove\_dyn} act on a whole vertex: its two bilinears
enter or leave the trace together, at the two times drawn by that one move. For a
spin-flip vertex $S^+(\t_1)S^-(\t_2)$, the vertex can only be removed when its own $S^+$
and $S^-$ are adjacent in the flip sequence, so the reachable configurations are those
that can be dismantled one adjacent or nested pair at a time. Pairings in which
\emph{every} vertex crosses another are unreachable, and the smallest of these has three
vertices --- so the two moves alone are not ergodic, and the omission is missing weight,
i.e.\ a bias, not merely slow mixing.

\texttt{swap\_dyn} (\texttt{moves/swap\_dyn.cpp}, after CT-SEG's
\texttt{swap\_spin\_lines}) supplies the missing connectivity. It picks two distinct
vertices and one bilinear in each, and exchanges their \emph{partners} --- which $S^+$ goes
with which $S^-$ --- by swapping the two bilinears' times between the vertices and
rebuilding each in catalog form ($O_1$ at the later time).

The exchange is restricted to two bilinears \emph{of the same kind}, so every fermion
operator stays at exactly the time it already occupied. The trace, the determinants, the
permutation sign and the Lang--Firsov weight (Eq.~\eqref{eq:ratio}, which depends only on
where the density kinks sit) are therefore all unchanged, only the coupling each pair is
evaluated at changes, and the acceptance is the ratio of coupling products,
\begin{equation}
  \frac{f_{t_1}(\t_1'-\t_2')\,f_{t_2}(\t_1''-\t_2'')}{f_{t_1}(\t_1-\t_2)\,f_{t_2}(\t_1''-\t_2'')} .
\end{equation}
No trace evaluation is needed. Two guards keep detailed balance, both symmetric between
the two directions: the ordered bilinear pair must match exactly one catalog entry (if it
matches several the vertex's entry is ambiguous and the move rejects in \emph{both}
directions), and the re-paired vertex must match an entry that exists.

Since no operator position changes, the move samples the weight of
Eq.~\eqref{eq:ssamp}: it adds connectivity, not a different measure. It is a prerequisite
for the histogram estimator below, which reads only vertex separations and so cannot see
an unreachable pairing at all.

\subsection{Stochastic channels: the vertex histogram}
\label{sec:histogram}

Perturb one type, $f_t\to f_t+\epsilon g$. In the sampled weight every vertex of type
$t$ carries a factor $-(f_t+\epsilon g)(s_v)$, so
$\partial_\epsilon\ln Z|_{0}=\aver{\sum_{v\in t}g(s_v)/f_t(s_v)}$. From
Eq.~\eqref{eq:ssamp}, with translation invariance,
$C_t(s)\equiv\aver{\TT\,O_{1,t}(s)\,O_{2,t}(0)}$ for $s\in(0,\beta)$, and
$\int_0^\beta d\t_2\int_{\t_2}^\beta d\t_1\,F(\t_1-\t_2)=\int_0^\beta ds\,(\beta-s)F(s)$,
we get $\partial_\epsilon\ln Z=-\int_0^\beta ds\,(\beta-s)\,g(s)\,C_t(s)$. Since this
holds for every $g$, the histogram of vertex separations is
\begin{equation}
  H_t(s)\equiv\aver{\sum_{v\in t}\delta(s-s_v)}=-(\beta-s)\,f_t(s)\,C_t(s).
  \label{eq:hist}
\end{equation}
The factor $(\beta-s)$ is the length of the ordered region at separation $s$: a
short-separation pair fits into $[0,\beta]$ in more ways than a long one.

\paragraph{Fold.} Let $\bar t$ be the reversed type ($O_{1,\bar t}=O_{2,t}$,
$O_{2,\bar t}=O_{1,t}$: $(b,a)$ for $(a,b)$, $S^-S^+$ for $S^+S^-$, and $\bar t=t$ for
$(a,a)$). By cyclicity of the trace, $C_{\bar t}(\beta-s)=C_t(s)$. Adding
Eq.~\eqref{eq:hist} for $t$ at $s$ and for $\bar t$ at $\beta-s$ gives
\begin{equation}
  C_t(s)=-\,\frac{H_t(s)+H_{\bar t}(\beta-s)}{(\beta-s)\,f_t(s)+s\,f_{\bar t}(\beta-s)},
  \label{eq:histest}
\end{equation}
with $f_{\bar t}=0$ if $\bar t$ is not in the catalogue. Under Eq.~\eqref{eq:dsym} the
denominator is $\beta f_t(s)$. The fold matters: without it the denominator
$(\beta-s)f_t(s)$ vanishes as $s\to\beta$.

\paragraph{Legendre form.} Accumulate $h_{t,l}=\aver{\sum_{v\in t}P_l(x(s_v))}$. Then
$H_t(s)=\sum_l\frac{2l+1}{\beta}h_{t,l}P_l(x(s))$, and since $P_l(-x)=(-1)^lP_l(x)$ the
folded numerator has coefficients $h_{t,l}+(-1)^lh_{\bar t,l}$. The division in
Eq.~\eqref{eq:histest} is done on the $\t$ grid.

\paragraph{Hellmann--Feynman check.} Integrating Eq.~\eqref{eq:hist} gives
$\aver{k_t}=-\int_0^\beta(\beta-s)f_t(s)C_t(s)\,ds=\lambda\partial_\lambda\ln Z$ for
$f_t\to\lambda f_t$, where $k_t$ is the number of type-$t$ vertices. This is a scalar
consistency check on the reconstructed $C_t$.

\paragraph{Properties.}
\begin{itemize}
  \item \emph{Exact}, because it is a derivative of the sampled weight, which equals $Z$
    only if the sampling is ergodic (\texttt{swap\_dyn}) and the vertex operators carry
    their Lang--Firsov dressing (\texttt{qmc\_data::trace\_ops}).
  \item \emph{Free}: it uses only the vertex times, with no extra trace evaluations.
  \item \emph{Noisy where $|f_t|$ is small}: the variance diverges as the coupling goes
    to zero, like the $M^{-1}$ estimator of $G$ as $\Delta\to0$. It only exists for
    channels that are actually sampled with $f_t\neq0$. Under the split of
    Sec.~\ref{sec:split} these are the residual $R$ vertices: for a density pair,
    Eq.~\eqref{eq:histest} gives the full-model $\aver{n_a(s)n_b(0)}$ wherever
    $R_{ab}\neq0$.
  \item A dynamical spin-flip channel gives $\aver{S^+(s)S^-(0)}$, which with SU(2)
    symmetry must equal $2\aver{S_z(s)S_z(0)}$ from the kink estimator. This is an
    exact cross-check between the two estimators.
\end{itemize}

\subsection{Lang--Firsov channels: the kink estimator is the same derivative}
\label{sec:kink-derivative}

Take a perturbation $E$ in the Lang--Firsov space $\mathcal L$ of Sec.~\ref{sec:split}
($PEP=E$), obeying Eq.~\eqref{eq:dsym}. The Lang--Firsov representation of
$Z[D+\epsilon E]$ is exact for every $\epsilon$, and linear in $E$ through two pieces:
the factor $\exp\{\sum_{\alpha<\beta}s_\alpha s_\beta K^E_{a_\alpha b_\beta}\}$ and the
static shift $-\tfrac\epsilon2\sum_{ab}K'^E_{ab}(0)\,n_an_b$ in $\Hloc$ (with
$n_an_a=n_a$). Hence
\[
  \partial_\epsilon\ln Z=\aver{\sum_{\alpha<\beta}s_\alpha s_\beta
  K^E_{a_\alpha b_\beta}(\t_\alpha-\t_\beta)}
  +\frac\beta2\sum_{ab}K'^E_{ab}(0)\aver{n_an_b}.
\]
On the physical side, with $E$ and $\chi_{ab}(s)=\aver{\TT\,n_a(s)n_b(0)}$ both periodic
on the full square (so there is no $(\beta-s)$ factor here),
$\partial_\epsilon\ln Z=-\tfrac12\int\!\!\int\sum_{ab}E_{ab}(\t-\t')\aver{n_a(\t)n_b(\t')}
=-\tfrac\beta2\sum_{ab}\int_0^\beta E_{ab}(s)\chi_{ab}(s)\,ds$.

Expand in Legendre polynomials with the normalization of Eq.~\eqref{eq:dn}:
$\int_0^\beta E\chi\,ds=\sum_n\frac{\beta}{2n+1}e_n\chi_n$; $K^E$ has coefficients
$k=Me$ (\texttt{build\_K\_n}); and $K'^E(0)=-\beta e_0$ by Eq.~\eqref{eq:kprime0}, since
$e_1=0$ under Eq.~\eqref{eq:dsym}. Define the ordered-pair kink correlator
\[
  \alpha^{ab}_p=\aver{\sum_{\alpha\neq\beta}s_\alpha s_\beta\,
  P_p\big(x(\t_\alpha-\t_\beta\bmod\beta)\big)},\qquad
  \alpha\ \text{on orbital}\ a,\ \beta\ \text{on orbital}\ b,
\]
which is exactly what \texttt{measure\_D0\_corr} accumulates, over all trace operators
including stochastic-vertex ones. Because $K^E_{ab}(t)=K^E_{ab}(\beta-t)=K^E_{ba}(t)$,
the folded sum over $\alpha<\beta$ equals half the ordered sum:
$\sum_{\alpha<\beta}s_\alpha s_\beta K^E=\tfrac12\sum_{ab}\sum_n e^{ab}_n(M^T\alpha^{ab})_n$.
Equating the coefficients of $e^{ab}_n$ gives
\begin{equation}
  \chi^{ab}_n=-\frac{2n+1}{\beta^2}\big(M^T\alpha^{ab}\big)_n
  \;+\;\aver{n_an_b}\,\delta_{n0}.
  \label{eq:kinkest}
\end{equation}
The first term is line for line the code's \verb|q_n = -sum_p M(p,n) alpha(p) (2n+1)/beta^2|.
The second term is the $\aver{n_an_b}$ offset that \texttt{solver.py} adds to
$Q_l[0]$. So the offset is not an ad hoc fix: it is the derivative of the $K'(0)$ static
shift. (Allowing $e_1\neq0$ for $a\neq b$ adds $-\tfrac\beta3e_1$ to $K'(0)$, which
cancels between $(a,b)$ and $(b,a)$ because $\aver{n_an_b}=\aver{n_bn_a}$.) Three
consequences:
\begin{itemize}
  \item \textbf{Only conserved components are determined.} $E\in\mathcal L$ means
    $e^{ab}_n=\sum_{ij}v_i[a]v_j[b]\gamma^{ij}_n$, so Eq.~\eqref{eq:kinkest} fixes only
    the contractions $\sum_{ab}v_i[a]v_j[b]\chi^{ab}_n=\aver{O_iO_j}_n$, and nothing
    else. If some $n_a$ is not conserved, the orbital-resolved $\chi_{ab}$ it would
    report is not a measurement of anything; the conserved-basis output is the complete
    and correct result.
  \item \textbf{No Lang--Firsov coupling is needed.} The identity holds at any base
    point, including zero Lang--Firsov coupling, because the representation of
    $Z[D+\epsilon E]$ is exact for all $\epsilon$. This is why \texttt{measure\_D0\_corr}
    does not require \texttt{lang\_firsov}.
  \item \textbf{The contact pairs drop out.} The two operators of one stochastic density
    vertex sit one tick apart and contribute $K^E(0^+)=0$ to the weight, so skipping
    them in $\alpha$ (as \texttt{measure\_D0\_corr} does) is exact, not an
    approximation.
\end{itemize}
Supporting evidence for the all-densities-conserved case: $\chi_{S_zS_z}$ from this
estimator agrees with CT-INT and CT-SEG for the spin--spin model. The ED reference with
a phonon and a discrete bath (to be added) tests the multi-orbital case independently.

\section{Tests, benchmarks, and open items}
\label{sec:tests}

\subsection{What is covered}

Registered \texttt{ctest}s, in the order they exercise the pipeline:
\begin{itemize}
  \item \texttt{Py\_dyn\_static\_shift} --- the static shift of
    Sec.~\ref{sec:kprime-practice} on its own, no Monte Carlo: pins
    \texttt{kprime\_0}/\texttt{static\_shift}/\texttt{half\_filling\_mu} to
    \texttt{ed\_reference}'s ED-validated $\mu$ and to \texttt{holstein.py}'s
    $U/2-g^2/\w_0^2$, and asserts the two orbitals of the $g=(0.7,0.3)$ case differ (the
    sign trap).
  \item \texttt{spin\_spin} (C++) and \texttt{Py\_spin\_spin} --- single orbital, the
    scalar \texttt{Jperp\_tau}/\texttt{D0\_tau} path.
  \item \texttt{Py\_kanamori\_dyn} --- a Kanamori $\Hloc$ with spin-flip and pair-hopping,
    so no individual orbital density commutes with it while $N_\uparrow,N_\downarrow$ do.
    Deterministic, 5000 cycles, against a reference file; all 16 vertices route
    analytically.
  \item \texttt{Py\_kanamori\_dyn\_nonuniform} --- the same with $U(\t)\neq U'(\t)$,
    confirming the claim at the end of Sec.~\ref{sec:conserved} that the standard Kanamori
    shape is exactly representable (16/16 analytic).
  \item \texttt{Py\_kanamori\_dyn\_selfconsistency} --- the physics check:
    \texttt{lang\_firsov=True} against forced \texttt{lang\_firsov=False} on an identical
    setup, compared through $G_l$. Agreement ${\sim}0.02$ at 200k/300k cycles, average
    sign ${\sim}0.9998$ in both.
\end{itemize}
Against other solvers, for the single-orbital case CT-SEG can reach: with
\texttt{lang\_firsov=True} and the same bath, both give average sign exactly $1.0$ and
$\max|G_{\rm ctseg}-G_{\rm cthyb}|$ falls from $3.5$ at 2k cycles to $0.17$ at 100k, which
validates the $D0$ sign convention and the Lang--Firsov path. $\chi_{S_zS_z}$ from the
estimator of Sec.~\ref{sec:derivative} agrees with CT-INT and CT-SEG for the spin--spin
model.

Production benchmarks under \texttt{benchmark/dynamic\_int/} are not CI-gated:
\texttt{kanamori\_dyn.py} (Bethe bath, with the $\mu$ correction applied from the input
coupling), \texttt{kanamori\_dyn\_spinflip.py} (a retarded inter-orbital spin-flip channel
on top of the density phonon --- mixed routing, 16 analytic and 4 stochastic for
$n_{\rm orb}=2$), \texttt{multiorb/} (solid\_dmft $GW$+EDMFT input), and
\texttt{dca\_spin\_spin/} (below). Each has its own launcher and plotting script.

\subsection{The DCA benchmark}

\texttt{benchmark/dynamic\_int/dca\_spin\_spin/} is a two-patch DCA calculation (VBDMFT
patches, $|k_x|,|k_y|<\pi/\sqrt2$ and the rest) with a \emph{real-space} retarded full
$\mathbf S\!\cdot\!\mathbf S$ interaction. The static $U$ and the retarded
$\mathbf S\!\cdot\!\mathbf S$ are both local on the two \emph{rotated} cluster sites, so in
the solver's working (patch) basis the dynamical vertices are off-diagonal bilinears
$c\ndag_{Ks}c_{K's'}$ with $K\neq K'$. The per-vertex test of Sec.~\ref{sec:eligibility}
rejects \emph{every} vertex here, since $n_{Ks}$ does not commute with $\Hloc$, so the
whole coupling is handled by the ${\rm span}\{N_\uparrow,N_\downarrow\}$ split of
Sec.~\ref{sec:split}. \texttt{check\_rotation.py} performs 22 symbolic checks with no
solver, among them the SU(2) consistency relation between the longitudinal and transverse
parts of a $\lambda(\t)\mathbf S\!\cdot\!\mathbf S$ coupling,
\texttt{Jperp\_tau}${}=4D0_{\rm same-spin}$ (same sign), which follows because
$\lambda S^zS^z$ gives $\lambda/4$ for a same-spin pair while
\texttt{expand\_Jperp\_into\_vertices} consumes $\texttt{Jperp\_tau}/2$ per ordering of
$\tfrac\lambda2(S^+S^-+S^-S^+)$. \texttt{expand\_retarded\_product} in \texttt{dca\_model.py}
expands the rotated-basis spin operators into monomial pairs and merges duplicates: for
$J_{\rm intra}=0,\ J_{\rm inter}=0.5$ that is 192 raw pairs down to 48 vertices (16
density-density, 32 off-diagonal); for uniform $J$, 384 down to 24, because
$\sum_{ij}\mathbf S_i\!\cdot\!\mathbf S_j=\mathbf S_{\rm tot}^2$ is basis independent and
the rotation cancels.

\subsection{Open items}

\begin{itemize}
  \item \textbf{An ED reference for the DCA benchmark} is not written yet, so that case is
    currently checked only for internal consistency.
  \item \textbf{A general sign-preserving split.} Eq.~\eqref{eq:signsplit} applies only
    when the conserved combinations are indicator vectors of disjoint orbital sets. For
    any other $\mathcal C$ the rejected density vertices stay fully stochastic; the general
    case is a small linear program per $\t$, not implemented.
  \item \textbf{The single-orbital pure-stochastic sign problem.} \texttt{spin\_spin.cpp}
    with \texttt{lang\_firsov=False} has an average sign of ${\sim}0.009$, effectively
    noise, against its reference file. Some sign problem is intrinsic here; this magnitude
    is not. The $D0$ sign convention is validated against CT-SEG and the scalar
    \texttt{Jperp\_tau} gives bit-identical results, so the open candidates are a stale
    reference file and a bug in the \texttt{insert\_dyn}/\texttt{remove\_dyn} weight. Two
    data points favour the former: \texttt{kanamori\_dyn\_selfconsistency} holds average
    sign near $1.0$ at 300k cycles with the same moves forced on, and the multi-orbital
    dynamical spin-flip channel degrades smoothly ($0.97/0.89/0.86/0.75$ at
    $g_{\rm sf}=0.3/0.5/0.7/1.0$) rather than collapsing.
  \item \textbf{\texttt{benchmark/dynamic\_int/CTHYB\_Benchmarking.ipynb}} parses, but its
    content has not been reviewed against the merge-resolution policy for this branch
    (default to \texttt{unstable}'s side unless there is a verified reason to prefer the
    other).
\end{itemize}

\subsection{Conventions}

\begin{itemize}
  \item \textbf{Infer less; require the caller to specify more.} An interface whose
    assumptions are not met raises an error instead of falling back quietly: this is why
    \texttt{expand\_Jperp\_into\_vertices} refuses an ambiguous block layout,
    \texttt{extract\_bilinear} refuses a non-unit coefficient, and a missing diagonal
    self-term is never supplied on the caller's behalf.
  \item \textbf{Nothing silent is approximate.} A coupling is resummed analytically only
    where that is provably exact; the remainder is separated by an identity, never by a
    heuristic that could quietly change the physics being asked for.
  \item \textbf{No unicode or mathematical symbols in identifiers} --- plain English names,
    however verbose.
\end{itemize}

\begin{thebibliography}{9}
\bibitem{LangFirsov1962} I.~G.~Lang and Yu.~A.~Firsov, \emph{Zh. Eksp. Teor. Fiz.}
  \textbf{43}, 1843 (1962).
\bibitem{WernerMillis2007} P.~Werner and A.~J.~Millis, \emph{Phys. Rev. Lett.}
  \textbf{99}, 146404 (2007).
\bibitem{WernerMillis2010} P.~Werner and A.~J.~Millis, \emph{Phys. Rev. Lett.}
  \textbf{104}, 146401 (2010).
\bibitem{GullRMP2011} E.~Gull, A.~J.~Millis, A.~I.~Lichtenstein, A.~N.~Rubtsov,
  M.~Troyer, and P.~Werner, \emph{Rev. Mod. Phys.} \textbf{83}, 349 (2011).
\bibitem{Nomura2015} Y.~Nomura, S.~Sakai, M.~Capone, and R.~Arita, \emph{Sci. Adv.}
  \textbf{1}, e1500568 (2015).
\bibitem{Steiner2015} K.~Steiner, Y.~Nomura, and P.~Werner, \emph{Phys. Rev. B}
  \textbf{92}, 115123 (2015).
\bibitem{Georges2013} A.~Georges, L.~de'~Medici, and J.~Mravlje, \emph{Annu. Rev.
  Condens. Matter Phys.} \textbf{4}, 137 (2013).
\end{thebibliography}

\end{document}
